% Solution to Problem 4 — The finite free Stam inequality (Srivastava)
% v8, corrected/completed against First Proof paper
% (Garza Vargas-Srivastava-Stier, "The finite free Stam inequality")

\documentclass[11pt]{article}

\usepackage[margin=1in]{geometry}
\usepackage{amsmath, amssymb, amsthm}
\usepackage{mathtools}
\usepackage{xcolor}
\usepackage{hyperref}
\hypersetup{colorlinks=true,
            linkcolor=blue!50!black,
            citecolor=blue!50!black,
            urlcolor=blue!50!black}

\newcommand{\R}{\mathbb{R}}
\newcommand{\Z}{\mathbb{Z}}
\newcommand{\bn}{\boxplus_n}
\newcommand{\score}{J_n}
\newcommand{\Var}{\operatorname{Var}}
\newcommand{\Hess}{\operatorname{Hess}}
\newcommand{\PSD}{\succeq 0}
\DeclareMathOperator{\sgn}{sgn}

\newtheorem{theorem}{Theorem}
\newtheorem{lemma}{Lemma}
\newtheorem{proposition}{Proposition}
\newtheorem{corollary}{Corollary}
\theoremstyle{definition}
\newtheorem{observation}{Observation}

\title{Solution to Problem 4: The Finite Free Stam Inequality}
\author{}
\date{May 27, 2026}

\begin{document}
\maketitle

\section*{Statement and Answer}

For a monic real-rooted polynomial $p$ of degree $n$ with roots
$\alpha=(\alpha_1,\dots,\alpha_n)$, define the \emph{finite free score vector}
$\score(\alpha)\in(\R\cup\{\infty\})^n$ by
\[
\score(\alpha)_i \;=\; \sum_{j\ne i}\frac{1}{\alpha_i-\alpha_j},
\]
and the \emph{finite free Fisher information}
$\Phi_n(p):=\|\score(\alpha)\|^2$, with $\Phi_n(p)=\infty$ if $p$ has a
multiple root. Let $\bn$ denote the finite free convolution of degree-$n$
polynomials.

\medskip
\noindent\textbf{Theorem.} \emph{For any two monic real-rooted polynomials
$p,q$ of degree $n$,}
\[
\frac{1}{\Phi_n(p\bn q)}\;\ge\;\frac{1}{\Phi_n(p)}+\frac{1}{\Phi_n(q)}.
\]

\medskip
\noindent The answer is \textbf{YES}. The result was conjectured by
D.~Shlyakhtenko; the proof below follows Blachman's strategy for the classical
Stam inequality, with the key analytic input replaced by a convexity statement
for hyperbolic polynomials.

\section*{Preliminaries on $\bn$}

We recall the standard description of the finite free convolution. If $\alpha$
and $\beta$ are vectors of roots for $p$ and $q$, then
\begin{equation}\label{eq:walsh}
p(x)\bn q(x)\;=\;\sum_{\pi\in S_n}\prod_{i=1}^n\bigl(x-\alpha_i-\beta_{\pi(i)}\bigr),
\end{equation}
and Walsh's theorem guarantees that $p\bn q$ is again real-rooted
\cite{Walsh, MSS22}. Thus $\bn$ induces a map
\[
\Omega_{\bn}:\R^n\times\R^n\to\R^n,
\]
where $\Omega_{\bn}(\alpha,\beta)$ is the ordered root vector of $p\bn q$. Write
$m_k(\alpha)=\frac1n\sum_i\alpha_i^k$ and
$\Var(\alpha)=m_2(\alpha)-m_1(\alpha)^2$, and let $\mathbf 1_n$ be the all-ones
vector.

\begin{proposition}[properties of $\bn$]\label{prop:props}
If $\gamma=\Omega_{\bn}(\alpha,\beta)$ then
\begin{enumerate}\setlength\itemsep{2pt}
\item (additivity) $m_1(\gamma)=m_1(\alpha)+m_1(\beta)$ and
$\Var(\gamma)=\Var(\alpha)+\Var(\beta)$;
\item (translation) $\Omega_{\bn}(\alpha+t\mathbf 1_n,\beta)=\gamma+t\mathbf 1_n$
and likewise in the second argument.
\end{enumerate}
\end{proposition}

\begin{proof}
(i) follows from \eqref{eq:walsh} together with the Newton identities applied
to the first two coefficients of $p\bn q$; (ii) is immediate from
\eqref{eq:walsh}.
\end{proof}

\section*{Score vectors via the heat flow}

The link between $\Phi_n$ and the dynamics of roots is the following, due to
Shlyakhtenko.

\begin{lemma}[score vector as a velocity]\label{lem:heat}
Assume $p$ has simple roots. Let $p_t(x):=\exp\!\bigl(-\tfrac{t}{2}\partial_x^2\bigr)p(x)$
and let $\alpha(t)$ be the ordered root vector of $p_t$. Then
$\alpha_i'(0)=\sum_{j\ne i}\frac{1}{\alpha_i-\alpha_j}$; that is,
$\alpha'(0)=\score(\alpha)$.
\end{lemma}

\begin{proof}
Roots stay simple near $t=0$. Implicitly differentiating
$p(\alpha_i(t))-\tfrac{t}{2}p''(\alpha_i(t))+t^2R(\alpha_i(t),t)=0$ at $t=0$
gives $\alpha_i'(0)=\tfrac12\,p''(\alpha_i)/p'(\alpha_i)$, which equals
$\sum_{j\ne i}(\alpha_i-\alpha_j)^{-1}$.
\end{proof}

The essential point is that $\exp(-\tfrac t2\partial_x^2)$ \emph{commutes with}
$\bn$ (finite free convolution commutes with any constant-coefficient
differential operator \cite{MSS22}). We exploit this twice.

\section*{Reduction to simple roots}

\begin{lemma}\label{lem:simple}
Let $T_\epsilon:=(1-\epsilon\,d/dx)^n$. For monic real-rooted $p$ of degree $n$:
(i) $T_\epsilon p$ is monic real-rooted with simple roots; (ii)
$\Phi_n(T_\epsilon p)\to\Phi_n(p)$ as $\epsilon\to0$; (iii)
$(T_\epsilon p)\bn(T_\epsilon q)=T_\epsilon^2(p\bn q)$.
\end{lemma}

\begin{proof}
(i) is Nuij's theorem \cite{Nuij}; (ii) holds since $\Phi_n$ is continuous in
the roots and the roots are continuous in $\epsilon$; (iii) holds because $\bn$
commutes with the differential operator $1-\epsilon\,d/dx$ \cite{MSS22}.
\end{proof}

By Lemma~\ref{lem:simple}(iii) and (ii), it suffices to prove the Theorem when
$p$, $q$, and $p\bn q$ all have simple roots; we assume this henceforth. Then
$\alpha,\beta,\gamma=\Omega_{\bn}(\alpha,\beta)$ have distinct entries and are
smooth functions of the coefficients of the corresponding polynomials. Let
$J_{\bn}$ denote the Jacobian of $\Omega_{\bn}$ at $(\alpha,\beta)$.

\section*{Step 1: score vectors and the Jacobian}

\begin{observation}\label{obs:scores}
For all $a,b\ge0$,
\[
J_{\bn}\bigl(a\,\score(\alpha),\,b\,\score(\beta)\bigr)\;=\;(a+b)\,\score(\gamma).
\]
\end{observation}

\begin{proof}
Set $p_t=\exp(-\tfrac t2\partial_x^2)p$ with ordered roots $\alpha(t)$, and
define $q_t,\beta(t)$ and $r_t,\gamma(t)$ analogously. Because $\bn$ commutes
with the heat operator,
\[
r_{(a+b)t}\;=\;p_{at}\bn q_{bt},
\]
so $\gamma((a+b)t)=\Omega_{\bn}(\alpha(at),\beta(bt))$. Differentiating at
$t=0$ and using the chain rule,
\[
(a+b)\,\gamma'(0)\;=\;J_{\bn}\!\begin{pmatrix} a\,\alpha'(0)\\ b\,\beta'(0)\end{pmatrix}.
\]
Lemma~\ref{lem:heat} identifies $\alpha'(0)=\score(\alpha)$, etc., giving the
claim.
\end{proof}

\section*{Step 2: the Jacobian is a contraction on mean-zero vectors}

This is the substance of the proof.

\begin{proposition}\label{prop:contraction}
If $u,v\in\R^n$ are each orthogonal to $\mathbf 1_n$, then
\[
\|J_{\bn}(u,v)\|^2\;\le\;\|u\|^2+\|v\|^2.
\]
\end{proposition}

Write $\gamma_i=\Omega_{\bn,i}(\alpha,\beta)$ for the coordinate functions and
$H^{(i)}:=\Hess\Omega_{\bn,i}$. Let
$V=\{(u,v)\in\R^n\times\R^n:u^\ast\mathbf 1_n=v^\ast\mathbf 1_n=0\}$.

\begin{lemma}[Hessian identity]\label{lem:hess}
For all $w\in V$,
\[
w^\ast J_{\bn}J_{\bn}^\ast w
\;=\;w^\ast\Bigl(I_n\oplus I_n-\sum_{i=1}^n\gamma_i H^{(i)}\Bigr)w.
\]
\end{lemma}

\begin{proof}
Fix $w=(u,v)\in V$ and set $\alpha(t)=\alpha+tu$, $\beta(t)=\beta+tv$,
$\gamma(t)=\Omega_{\bn}(\alpha(t),\beta(t))$. By the additivity of variance
(Proposition~\ref{prop:props}(i)),
\[
m_2(\gamma(t))-m_1(\gamma(t))^2
=m_2(\alpha(t))+m_2(\beta(t))-m_1(\alpha(t))^2-m_1(\beta(t))^2 .
\]
Since $(u,v)\in V$, the means $m_1(\alpha(t)),m_1(\beta(t))$ — hence
$m_1(\gamma(t))$ — are constant in $t$. Differentiating twice at $t=0$,
\[
\partial_t^2 m_2(\gamma(t))\big|_0=\partial_t^2\bigl(m_2(\alpha(t))+m_2(\beta(t))\bigr)\big|_0 .
\]
The left side equals
$\tfrac1n\sum_i\partial_t^2(\gamma_i^2)
=\tfrac2n\sum_i\bigl((\partial_t\gamma_i)^2+\gamma_i\partial_t^2\gamma_i\bigr)
=\tfrac2n\bigl(w^\ast J_{\bn}J_{\bn}^\ast w+\sum_i\gamma_i w^\ast H^{(i)}w\bigr)$,
while the right side equals $\tfrac2n\,w^\ast w$. Rearranging gives the
identity.
\end{proof}

The minus sign means we must show $\sum_i\gamma_i H^{(i)}\PSD$ on $V$. This is
where real-rootedness of $\bn$ is used in an essential way, via the theory of
hyperbolic polynomials.

\begin{theorem}[Bauschke--G\"uler--Lewis--Sendov \cite{BGLS}]\label{thm:bgls}
Let $f\in\R[x_1,\dots,x_m]$ be hyperbolic in the direction $w\in\R^m$, and for
$a\in\R^m$ let $\lambda_1(a)\ge\dots\ge\lambda_m(a)$ be the roots of
$t\mapsto f(a+tw)$. Then for each $k$, $\sigma_k(a):=\sum_{i=1}^k\lambda_i(a)$ is
convex in $a$.
\end{theorem}

\begin{corollary}\label{cor:psd}
For any reals $c_1\ge\dots\ge c_n$, the matrix $\sum_{i=1}^n c_i H^{(i)}$ is PSD.
\end{corollary}

\begin{proof}
Consider $f(x,a_1,\dots,a_n,b_1,\dots,b_n)=\sum_{\pi\in S_n}\prod_{i=1}^n
(x-a_i-b_{\pi(i)})$. By \eqref{eq:walsh} and Walsh's theorem $f$ is homogeneous
and hyperbolic in the direction $e_1=(1,0,\dots,0)$. By Theorem~\ref{thm:bgls}
each $\sigma_k(x,a,b)=\sum_{i=1}^k\lambda_i(x,a,b)$ is convex, where
$\lambda_i(x,a,b)$ are the roots of $f((x,a,b)+te_1)$. Since the $c_i$ are
ordered,
\[
L(x,a,b):=\sum_{i=1}^n c_i\lambda_i(x,a,b)
\]
is a nonnegative combination of the $\sigma_k$ (by Abel summation), hence convex;
so $\Hess L=\sum_i c_i\Hess\lambda_i\PSD$ everywhere. Evaluating at $x=0$,
$a=\alpha$, $b=\beta$, we have $\Hess\lambda_i=H^{(i)}$, giving
$\sum_i c_i H^{(i)}\PSD$.
\end{proof}

\begin{proof}[Proof of Proposition~\ref{prop:contraction}]
Let $(u,v)\in V$. By Lemma~\ref{lem:hess},
\[
\|J_{\bn}(u,v)\|^2=(u,v)^\ast J_{\bn}J_{\bn}^\ast(u,v)
=\|u\|^2+\|v\|^2-\sum_{i=1}^n\gamma_i\,(u,v)^\ast H^{(i)}(u,v).
\]
The roots $\gamma_1\ge\dots\ge\gamma_n$ are ordered, so
Corollary~\ref{cor:psd} (with $c_i=\gamma_i$) gives $\sum_i\gamma_i H^{(i)}\PSD$,
whence $\sum_i\gamma_i(u,v)^\ast H^{(i)}(u,v)\ge0$. The claim follows.
\end{proof}

\section*{Step 3: Blachman's argument}

\begin{proof}[Proof of the Theorem]
First, $\score(\alpha)\perp\mathbf 1_n$: indeed
$\sum_i\sum_{j\ne i}(\alpha_i-\alpha_j)^{-1}=0$ since each unordered pair
contributes $+$ and $-$ once. Likewise $\score(\beta)\perp\mathbf 1_n$. So
Proposition~\ref{prop:contraction} applies to
$(a\,\score(\alpha),\,b\,\score(\beta))$:
\[
\bigl\|J_{\bn}\bigl(a\,\score(\alpha),b\,\score(\beta)\bigr)\bigr\|^2
\;\le\;a^2\|\score(\alpha)\|^2+b^2\|\score(\beta)\|^2 .
\]
Combining with Observation~\ref{obs:scores},
\[
(a+b)^2\,\Phi_n(p\bn q)\;\le\;a^2\,\Phi_n(p)+b^2\,\Phi_n(q),
\]
valid for all $a,b\ge0$. Choosing $a=1/\Phi_n(p)$ and $b=1/\Phi_n(q)$, the right
side becomes $1/\Phi_n(p)+1/\Phi_n(q)$ and the factor $(a+b)^2$ becomes
$\bigl(1/\Phi_n(p)+1/\Phi_n(q)\bigr)^2$, so
\[
\Bigl(\tfrac{1}{\Phi_n(p)}+\tfrac{1}{\Phi_n(q)}\Bigr)^2\Phi_n(p\bn q)
\;\le\;\tfrac{1}{\Phi_n(p)}+\tfrac{1}{\Phi_n(q)} .
\]
Dividing through yields
$\dfrac{1}{\Phi_n(p\bn q)}\ge\dfrac{1}{\Phi_n(p)}+\dfrac{1}{\Phi_n(q)}$.
The simple-root reduction (Lemma~\ref{lem:simple}) extends this to all
real-rooted $p,q$.
\end{proof}

\section*{Remarks}

The proof mirrors Blachman's \cite{Blachman} derivation of the classical Stam
inequality \cite{Stam}: there the score of a sum is a conditional expectation of
the marginal score and the inequality reduces to Cauchy--Schwarz. Here the role
of that conditional-expectation step is played by
Observation~\ref{obs:scores} (the score of $p\bn q$ is the image of the marginal
scores under the Jacobian $J_{\bn}$) together with the operator-norm bound
Proposition~\ref{prop:contraction}. Real-rootedness of $\bn$ is used exactly
once and essentially — through the hyperbolicity of \eqref{eq:walsh} and the
BGLS convexity theorem — which is why the inequality fails for convolutions that
do not preserve real-rootedness.

\begin{thebibliography}{9}
\bibitem{MSS22} A.~W.~Marcus, D.~A.~Spielman, N.~Srivastava.
  \emph{Finite free convolutions of polynomials}.
  Probab.\ Theory Related Fields \textbf{182} (2022), 807--848.
\bibitem{Walsh} J.~L.~Walsh.
  \emph{On the location of the roots of certain types of polynomials}.
  Trans.\ Amer.\ Math.\ Soc.\ \textbf{24} (1922), 163--180.
\bibitem{Nuij} W.~Nuij.
  \emph{A note on hyperbolic polynomials}.
  Math.\ Scand.\ \textbf{23} (1968), 69--72.
\bibitem{Blachman} N.~Blachman.
  \emph{The convolution inequality for entropy powers}.
  IEEE Trans.\ Inform.\ Theory \textbf{11} (1965), 267--271.
\bibitem{BGLS} H.~H.~Bauschke, O.~G\"uler, A.~S.~Lewis, H.~S.~Sendov.
  \emph{Hyperbolic polynomials and convex analysis}.
  Canad.\ J.\ Math.\ \textbf{53} (2001), 470--488.
\bibitem{Stam} A.~Stam.
  \emph{Some inequalities satisfied by the quantities of information of Fisher
  and Shannon}. Inform.\ Control \textbf{2} (1959), 101--112.
\end{thebibliography}

\end{document}